\subsubsection{Software Standards}

To ensure cross-platform compatibility, developers rely on unified industry standards. Currently, there are two major public standards, OpenXR and WebXR.


OpenXR is the industry-standard API for XR applications. It is the interface between an application and an in-process or out-of-process XR runtime system. The runtime system handles functionality such as frame composition, peripheral management, and raw tracking information \parencite{openxr2024spec}. Major game engines and the Meta Spatial SDK are built upon OpenXR to ensure hardware-agnostic development. 

A primary bottleneck in XR development is the deployment loop. Compiling a standalone APK and pushing it to a device can take 10~15 minutes per iteration. There are two solutions to avoid the compiling: (1) Link Mode: Developers often use the Oculus Link App to preview scenes in PC VR, though this requires constantly putting on and taking off the headset. (2) Simulators: The Meta XR Simulator is increasingly becoming the preferred method for rapid iteration, allowing developers to simulate head and controller movements directly on their PC without needing physical hardware.

The WebXR Device API provides the interfaces necessary to enable developers to build compelling, comfortable, and safe immersive applications on the web across a wide variety of hardware form factors \parencite{w3c2022webxr}.

WebXR serves as a higher-level abstraction that runs in the browser, combining older specifications from WebVR and WebAR. The WebXR standard shares core features with and can be viewed as a subset of OpenXR. While traditional game engines offer WebXR export support, there are fewer dedicated development stacks built exclusively for it. 

For security reasons, WebXR APIs are restricted to HTTPS.Consequently, developers must use port-forwarding technologies during local development. Thanks to the efficiency of the web stack, a WebXR application typically only requires a few seconds to bundle. It is worth noting, however, that web rendering technology is currently transitioning from WebGL2 to WebGPU, which may introduce API changes in the WebXR ecosystem.

Although web technologies are widely used across various domains, WebXR remains an emerging technology and is still considered relatively unstable. Therefore, this study adopts the OpenXR standard to ensure greater stability and cross-platform compatibility.